% ============================================================
%  Глава 5 — SLAM: вероятностная карта занятости
%  Файл: pi_nodes/nodes/slam_map_node.py
% ============================================================
\chapter{SLAM: вероятностная карта занятости на log-odds}
\label{ch:slam_map}

\begin{filebox}[title={Файл исходного кода}]
\filepath{pi\_nodes/nodes/slam\_map\_node.py} \quad (Python)
\end{filebox}

\section{Вероятностная модель карты}

Карта занятости представляет каждую ячейку сетки как бинарную случайную
величину: свободно ($\mathtt{F}$) или занято ($\mathtt{O}$).

\subsection{Log-odds представление}

Вместо хранения вероятности $P(\mathtt{O})$ используется \textbf{логарифм
отношения правдоподобий} (log-odds), устойчивый к крайним значениям $0$ и $1$:

\begin{equation}
  l = \log\frac{P(\mathtt{O})}{1 - P(\mathtt{O})} = \log\frac{P(\mathtt{O})}{P(\mathtt{F})}
  \label{eq:log_odds_def}
\end{equation}

Обратное преобразование:

\begin{equation}
  P(\mathtt{O}) = \frac{e^l}{1 + e^l} = \sigma(l)
  \label{eq:log_odds_inv}
\end{equation}

\textbf{Преимущества}: обновление --- простое сложение; нет умножения вероятностей.

\begin{filebox}[title={\filepath{slam\_map\_node.py}, строки 42--48 --- константы log-odds}]
\begin{lstlisting}[style=pythonstyle, numbers=left, firstnumber=42]
L_OCC  =  0.85   # log-odds increment for occupied cell
L_FREE = -0.40   # log-odds increment for free cell
L_MIN  = -2.0    # minimum log-odds (very confident free)
L_MAX  =  3.5    # maximum log-odds (very confident occupied)
L_THRESH_OCC  =  0.6   # above this -> occupied (output)
L_THRESH_FREE = -0.5   # below this -> free (output)
\end{lstlisting}
\end{filebox}

\textbf{Физический смысл значений log-odds:}

\begin{center}
\begin{tabular}{lll}
\toprule
\textbf{log-odds $l$} & \textbf{$P(\mathtt{O})$} & \textbf{Интерпретация}\\
\midrule
$-2.0$ & $11.9\%$ & Очень уверенно свободно\\
$-0.5$ & $37.8\%$ & Скорее свободно\\
$0.0$  & $50.0\%$ & Неизвестно\\
$+0.6$ & $64.6\%$ & Скорее занято\\
$+3.5$ & $97.1\%$ & Очень уверенно занято\\
\bottomrule
\end{tabular}
\end{center}

\section{Трассировка лучей (Ray Tracing)}

\begin{filebox}[title={\filepath{slam\_map\_node.py}, строки 256--285 --- \_trace\_ray()}]
\begin{lstlisting}[style=pythonstyle, numbers=left, firstnumber=256]
def _trace_ray(self, ox, oy, angle, range_m):
    cos_a = math.cos(angle)
    sin_a = math.sin(angle)

    step = CELL_SIZE_M * 0.7   # 3.5 cm step
    dist = 0.0
    hit_range = min(range_m, SENSOR_MAX_RANGE)

    # Mark cells along ray as FREE
    while dist < hit_range - OBSTACLE_THICKNESS:
        wx = ox + cos_a * dist
        wy = oy + sin_a * dist
        ci, cj = self._world_to_cell(wx, wy)
        if 0 <= ci < MAP_CELLS and 0 <= cj < MAP_CELLS:
            idx = cj * MAP_CELLS + ci
            self._grid[idx] = max(L_MIN, self._grid[idx] + L_FREE)
        dist += step

    # Mark obstacle cells as OCCUPIED
    if range_m < SENSOR_MAX_RANGE - 0.05:
        for d in range(3):
            dd = hit_range + d * CELL_SIZE_M * 0.5
            wx = ox + cos_a * dd
            wy = oy + sin_a * dd
            ci, cj = self._world_to_cell(wx, wy)
            if 0 <= ci < MAP_CELLS and 0 <= cj < MAP_CELLS:
                idx = cj * MAP_CELLS + ci
                self._grid[idx] = min(L_MAX, self._grid[idx] + L_OCC)
\end{lstlisting}
\end{filebox}

\textbf{Параметрическое уравнение луча:}

\begin{equation}
  \begin{pmatrix} w_x(t) \\ w_y(t) \end{pmatrix}
  = \begin{pmatrix} o_x \\ o_y \end{pmatrix}
  + t \begin{pmatrix} \cos\alpha \\ \sin\alpha \end{pmatrix}, \quad t \in [0, r]
  \label{eq:ray_parametric}
\end{equation}

\textbf{Байесовское обновление} (вычитание для свободных, прибавление для занятых):

\begin{equation}
  \boxed{l_{\text{new}} = \mathrm{clamp}\!\left(l + \Delta l,\; L_{\min},\; L_{\max}\right)}
  \label{eq:log_odds_update}
\end{equation}

где:
\[
  \Delta l = \begin{cases}
    L_{\text{FREE}} = -0.40 & \text{луч прошёл через ячейку}\\
    L_{\text{OCC}} = +0.85 & \text{луч достиг препятствия}
  \end{cases}
\]

\section{Ультразвуковой конус}

\begin{filebox}[title={\filepath{slam\_map\_node.py}, строки 247--254 --- сканирование конуса}]
\begin{lstlisting}[style=pythonstyle, numbers=left, firstnumber=247]
# Ultrasonic cone: trace multiple rays within the FOV
angle_start = robot_theta - SENSOR_FOV_RAD   # theta - 15 deg
angle_end   = robot_theta + SENSOR_FOV_RAD   # theta + 15 deg

angle = angle_start
while angle <= angle_end:
    self._trace_ray(robot_x, robot_y, angle, range_m)
    angle += ANGULAR_RESOLUTION              # step = 0.05 rad ~ 3 deg
\end{lstlisting}
\end{filebox}

Количество лучей на одно измерение:

\begin{equation}
  N_{\text{rays}} = \left\lfloor \frac{2 \cdot \text{SENSOR\_FOV\_RAD}}{\text{ANGULAR\_RESOLUTION}} \right\rfloor
  = \left\lfloor \frac{2 \times 0.26}{0.05} \right\rfloor = 10\,\text{ лучей}
  \label{eq:ray_count}
\end{equation}

\section{Преобразование координат}

\begin{filebox}[title={\filepath{slam\_map\_node.py}, строки 288--298 --- world $\leftrightarrow$ grid}]
\begin{lstlisting}[style=pythonstyle, numbers=left, firstnumber=288]
def _world_to_cell(self, wx, wy):
    ci = int((wx - self._origin_x) / CELL_SIZE_M)  # column
    cj = int((wy - self._origin_y) / CELL_SIZE_M)  # row
    return ci, cj

def _cell_to_world(self, ci, cj):
    wx = self._origin_x + (ci + 0.5) * CELL_SIZE_M
    wy = self._origin_y + (cj + 0.5) * CELL_SIZE_M
    return wx, wy
\end{lstlisting}
\end{filebox}

\textbf{Мировые координаты $\to$ ячейка сетки:}

\begin{equation}
  c_i = \left\lfloor \frac{w_x - o_x}{\delta} \right\rfloor, \quad
  c_j = \left\lfloor \frac{w_y - o_y}{\delta} \right\rfloor
  \label{eq:world_to_grid}
\end{equation}

\textbf{Ячейка $\to$ центр в мировых координатах:}

\begin{equation}
  w_x = o_x + (c_i + 0.5)\,\delta, \quad
  w_y = o_y + (c_j + 0.5)\,\delta
  \label{eq:grid_to_world}
\end{equation}

где $o_x = o_y = -5\,\text{м}$ (начало координат карты), $\delta = 0.05\,\text{м/ячейку}$.

\section{Кластеризация объектов с Евклидовой метрикой}

\begin{filebox}[title={\filepath{slam\_map\_node.py}, строки 170--188 --- \_detections\_cb()}]
\begin{lstlisting}[style=pythonstyle, numbers=left, firstnumber=172]
# Search for nearest existing object (same class+colour)
for obj_id, obj in self._objects.items():
    dx = obj['x'] - wx
    dy = obj['y'] - wy
    if math.sqrt(dx*dx + dy*dy) < OBJ_CLUSTER_RADIUS_M:  # 0.30 m
        matched_id = obj_id
        break

# Update position: sliding average
if matched_id:
    alpha = 0.3
    obj['x'] = round(alpha * wx + (1 - alpha) * obj['x'], 3)
    obj['y'] = round(alpha * wy + (1 - alpha) * obj['y'], 3)
\end{lstlisting}
\end{filebox}

Проверка принадлежности к кластеру:

\begin{equation}
  d_{\text{Euclid}} = \sqrt{(x_{\text{new}} - x_{\text{obj}})^2 + (y_{\text{new}} - y_{\text{obj}})^2}
  < 0.30\,\text{м}
  \label{eq:cluster_dist}
\end{equation}

Скользящее среднее позиции (EMA с $\alpha = 0.3$):

\begin{equation}
  \vect{p}_{\text{obj}} \leftarrow 0.3\,\vect{p}_{\text{new}} + 0.7\,\vect{p}_{\text{obj}}
  \label{eq:ema_position}
\end{equation}

\section{Параметры сетки}

\begin{center}
\begin{tabular}{lll}
\toprule
\textbf{Параметр} & \textbf{Значение} & \textbf{Описание}\\
\midrule
\texttt{MAP\_SIZE\_M}      & $10\,\text{м}$ & Размер карты ($\pm 5\,\text{м}$)\\
\texttt{CELL\_SIZE\_M}     & $0.05\,\text{м}$ & Размер ячейки (5 см)\\
\texttt{MAP\_CELLS}        & $200 \times 200$ & Размер сетки\\
\texttt{SENSOR\_FOV\_RAD}  & $0.26\,\text{рад}$ & Полу-угол конуса (${\approx}15^\circ$)\\
\texttt{SENSOR\_MAX\_RANGE}& $2.0\,\text{м}$ & Макс. дальность HC-SR04\\
\bottomrule
\end{tabular}
\end{center}
